\documentclass[12pt,a4paper]{ctexart}
\usepackage{amsmath}
\usepackage{amssymb}
\usepackage{booktabs}
\usepackage{float}
\usepackage{geometry}
\usepackage{graphicx}
\usepackage{caption}

\geometry{margin=2.5cm}
\captionsetup{font=small}
\graphicspath{{figures/}}

\title{实验名称}
\author{姓名，院系/班级，学号，组号，实验日期}
\date{}

\begin{document}

\maketitle

\begin{abstract}
% 用“目的-方法-结果-结论”组织摘要，可附关键词。
\end{abstract}

\section{实验原理}

% 结合必要公式和示意图说明实验原理。

\section{实验装置与实验方法}

% 说明仪器、装置和操作方法。
% 如有单独的装置图、电路图、波形图，建议放在 figures/ 目录并用相对路径插入。
% 如图像来自原始照片中的局部区域，应先裁剪为单独图片再插入。

\section{数据处理}

% 建议将数据表写入 data.tex 并在此处 \input{data.tex}。
% 计算建议写成：物理量 = 数据代入式 = 结果（单位）。
\input{data.tex}

\section{总结}

% 据实分析结果、误差来源、实验现象和改进建议。

\section{参考文献}

% 按课程要求列出参考文献。

\section{原始数据}

% 如课程要求，请插入带教师签字的原始数据照片或扫描件。
% 为保证 Overleaf 可直接编译，图片文件应保存在 figures/ 目录下。

\end{document}
